\chapter{2006年山东卷（文）}
\section{单选题}
\begin{problem}
定义集合运算： \( A \odot B = \{z \mid z = xy(x + y), x \in A, y \in B\} \) ，设集合 \( A = \{0, 1\} \) ， \( B = \{2, 3\} \) ，则集合 \( A \odot B \) 的所有元素之和为（\quad）

\choices
    {0}
    {6}
    {12}
    {18}
\end{problem}

\begin{problem}
设 \( f(x) = \begin{cases} 2\e^{x - 1}, & x < 2, \\ \log_3(x^2 - 1), & x \geqslant 2, \end{cases} \) 则 \( f[f(2)] \) 的值为（\quad）

\choices
    {0}
    {1}
    {2}
    {3}
\end{problem}

\begin{problem}
函数 \( y = 1 + a^x \) (\( 0 < a < 1 \)) 的反函数的图象大致是（\quad）

\begin{center}
\bitmapinclude[max width=0.72\linewidth,max totalheight=9.0cm]{img/2006/shandong_liberal/q3_options.png}
\end{center}
\end{problem}

\begin{problem}
设向量 \(a = (1, -3)\)，\(b = (-2, 4)\)，若表示向量 \(4a, 3b - 2a, c\) 的有向线段首尾相接能构成三边线，则向量 \(c\) 为（\quad）

\choices
    {\((1, - 1)\)}
    {\((-1, 1)\)}
    {\((-4, 6)\)}
    {\((4, - 6)\)}
\end{problem}

\begin{problem}
已知定义在 \(\R\) 上的奇函数 \(f(x)\) 满足 \(f(x + 2) = -f(x)\)，则 \(f(6)\) 的值为（\quad）

\choices
    {\(-1\)}
    {0}
    {1}
    {2}
\end{problem}

\begin{problem}
在 \(\triangle ABC\) 中，角 \(A, B, C\) 的对边分别为 \(a, b, c\)，已知 \(A = \frac{\pi}{3}, a = \sqrt{3}, b = 1\)，则 \(c =\)（\quad）

\choices
    {1}
    {2}
    {\(\sqrt{3} - 1\)}
    {\(\sqrt{3}\)}
\end{problem}

\begin{problem}
在给定双曲线中，过焦点且垂直于实轴的弦长为 \(\sqrt{2}\)，焦点到相应准线的距离为 \(\frac{1}{2}\)，则该双曲线的离心率为（\quad）

\choices
    {\(\frac{\sqrt{2}}{2}\)}
    {2}
    {\(\sqrt{2}\)}
    {\(2 \sqrt{2}\)}
\end{problem}

\begin{problem}
正方体的内切球与其外接球的体积之比为（\quad）

\choices
    {\( 1 : \sqrt{3} \)}
    {1:3}
    {\(1:3\sqrt{3}\)}
    {1:9}
\end{problem}

\begin{problem}
设 \( p: x^2 - x - 2 < 0 \)，\( q: \frac{1 + x}{|x| - 2} < 0 \)，则 \( p \) 是 \( q \) 的（\quad）

\choices
    {充分不必要条件}
    {必要不充分条件}
    {充要条件}
    {既不充分也不必要条件}
\end{problem}

\begin{problem}
已知 \(\left(x^{2} - \frac{1}{\sqrt[3]{x}}\right)^{n}\) 的展开式中第三项与第五项的系数之比为 \(\frac{3}{14}\)，则展开式中常数项是（\quad）

\choices
    {\(-1\)}
    {1}
    {\(-45\)}
    {45}
\end{problem}

\begin{problem}
已知集合 \(A = \{5\}\)，\(B = \{1, 2\}\)，\(C = \{1, 3, 4\}\)，从这三个集合各取一个元素构成空间直角坐标系中点的坐标，则确定的不同点的个数为（\quad）

\choices
    {33}
    {34}
    {35}
    {36}
\end{problem}

\begin{problem}
已知 \(x\) 和 \(y\) 是正整数，且满足约束条件 \(\left\{ \begin{array}{l} x + y \leqslant 10, \\ x - y \leqslant 2, \\ 2x \geqslant 7, \end{array} \right.\) 则 \(z = 2x + 3y\) 的最小值是（\quad）

\choices
    {24}
    {14}
    {13}
    {11.5}
\end{problem}

\section{填空题}
\begin{problem}
某学校共有师生 2400 人，现用分层抽样的方法，从所有师生中抽取一个容量为 160 的样本. 已知从学生中抽取的人数为 150，那么该学校的教师人数是\fillinblank{}.
\end{problem}

\begin{problem}
设 \(S_{n}\) 为等差数列 \(\{a_{n}\}\) 的前 \(n\) 项和，\(S_{4} = 14, S_{10} - S_{7} = 30\)，则 \(S_{9} = \)\fillinblank{}.
\end{problem}

\begin{problem}
已知抛物线 \( y^{2} = 4x \)，过点 \( P(4,0) \) 的直线与抛物线相交于 \( A(x_{1},y_{1}) \)，\( B(x_{2},y_{2}) \) 两点，则 \( y_{1}^{2} + y_{2}^{2} \) 的最小值是\fillinblank{}.
\end{problem}

\begin{problem}
如图，在正三棱柱 \(ABC - A_{1}B_{1}C_{1}\) 中，所有棱长均为 1 ，则点 \(B_{1}\) 到平面 \(ABC_{1}\) 的距离为\fillinblank{}.
\begin{center}
\bitmapinclude[max width=0.42\linewidth,max totalheight=4.8cm]{img/2006/shandong_liberal/q16_fig1.png}
\end{center}
\end{problem}

\section{解答题}
\begin{problem}
设函数 \( f(x) = 2x^{3} - 3(a - 1)x^{2} + 1 \)，其中 \( a \geqslant 1 \).

\begin{enumerate}
\item 求 \( f(x) \) 的单调区间；
\item 讨论 \( f(x) \) 的极值.
\end{enumerate}
\end{problem}

\begin{problem}
已知函数 \( f(x) = A \sin^2 (\omega x + \varphi) \left( A > 0, \omega > 0, 0 < \varphi < \frac{\pi}{2} \right) \)，且 \( y = f(x) \) 的最大值为2，其图象相邻两对称轴间的距离为2，并过点 \((1,2)\).

\begin{enumerate}
\item 求 \( \varphi \)；
\item 计算 \( f(1) + f(2) + \cdots + f(2008) \).
\end{enumerate}
\end{problem}

\begin{problem}
盒中装有标有数字 1, 2, 3, 4 的卡片各 2 张，从盒中任意抽取 3 张，每张卡片被取出的可能性都相等，求:

\begin{enumerate}
\item 抽出的 3 张卡片上最大的数字是 4 的概率；
\item 抽出的 3 张中有 2 张卡片上的数字是 3 的概率；
\item 抽出的 3 张卡片上的数字互不相同的概率.
\end{enumerate}
\end{problem}

\begin{problem}
如图，已知四棱锥 \(P - ABCD\) 的底面 \(ABCD\) 为等腰梯形，\(AB // DC\)，\(AC \perp BD\)，\(AC\) 与 \(BD\) 相交于点 \(O\)，且顶点 \(P\) 在底面上的射影恰为 \(O\) 点，又 \(BO = 2\)，\(PO = \sqrt{2}\)，\(PB \perp PD\).

\begin{enumerate}
\item 求异面直线 \(PD\) 与 \(BC\) 所成角的余弦值；
\item 求二面角 \(P - AB - C\) 的大小；
\item 设点 \(M\) 在棱 \(PC\) 上，且 \(\frac{PM}{MC} = \lambda\)，问 \(\lambda\) 为何值时，\(PC \perp\) 平面 \(BMD\).
\end{enumerate}

\begin{center}
\bitmapinclude[max width=0.42\linewidth,max totalheight=4.8cm]{img/2006/shandong_liberal/q20_fig1.png}
\end{center}
\end{problem}

\begin{problem}
已知椭圆的中心在坐标原点 \(O\)，焦点在 \(x\) 轴上，椭圆的短轴端点和焦点所组成的四边形为正方形，两准线间的距离为4.

\begin{enumerate}
\item 求椭圆的方程；
\item 直线 \(l\) 过点 \(P(0,2)\) 且与椭圆相交于 \(A\)、\(B\) 两点，当 \(\triangle AOB\) 面积取得最大值时，求直线 \(l\) 的方程.
\end{enumerate}
\end{problem}

\begin{problem}
已知数列 \(\{a_{n}\}\) 中，\(a_{1} = \frac{1}{2}\)，点 \((n, 2a_{n+1} - a_{n})\) 在直线 \(y = x\) 上，其中 \(n = 1, 2, 3, \dots\).

\begin{enumerate}
\item 令 \( b_{n} = a_{n+1} - a_{n} - 1 \)，求证数列 \( \{b_{n}\} \) 是等比数列；
\item 求数列 \(\{a_{n}\}\) 的通项；
\item 设 \(S_{n}\)、\(T_{n}\) 分别为数列 \(\{a_{n}\}\)、\(\{b_{n}\}\) 的前 \(n\) 项和.是否存在实数 \(\lambda\)，使得数列 \(\left\{\frac{S_{n} + \lambda T_{n}}{n}\right\}\) 为等差数列？若存在，试求出 \(\lambda\)；若不存在，则说明理由.
\end{enumerate}
\end{problem}
